\documentclass[11pt,letterpaper]{article}
\usepackage[T1]{fontenc}
\usepackage{lmodern}
\usepackage[margin=0.95in,headheight=15pt]{geometry}
\usepackage{amsmath,amssymb,amsthm,mathtools,booktabs,array}
\usepackage{mathrsfs}
\usepackage{microtype}
\usepackage{enumitem}
\usepackage{xurl}
\usepackage[hidelinks]{hyperref}
\usepackage{fancyhdr}
\pagestyle{fancy}\fancyhf{}
\fancyhead[L]{\small Continuous-depth Lerch calculus}
\fancyhead[R]{\small ProveIt continuation}
\fancyfoot[C]{\thepage}
\setlength{\parskip}{3pt}
\setlength{\parindent}{15pt}
\setlist{itemsep=3pt,topsep=5pt}
\allowdisplaybreaks[2]
\numberwithin{equation}{section}
\newtheorem{theorem}{Theorem}[section]
\newtheorem{lemma}[theorem]{Lemma}
\newtheorem{corollary}[theorem]{Corollary}
\newtheorem{proposition}[theorem]{Proposition}
\theoremstyle{definition}\newtheorem{definition}[theorem]{Definition}
\theoremstyle{remark}\newtheorem{remark}[theorem]{Remark}
\newcommand{\C}{\mathbb C}\newcommand{\R}{\mathbb R}
\newcommand{\N}{\mathbb N}\newcommand{\Q}{\mathbb Q}
\newcommand{\Rea}{\operatorname{Re}}\newcommand{\Ima}{\operatorname{Im}}
\newcommand{\Li}{\operatorname{Li}}\newcommand{\dd}{\,\mathrm d}
\newcommand{\F}{\mathcal F}\newcommand{\E}{\mathcal E}
\newcommand{\B}{\mathcal B}\newcommand{\PP}{\mathcal P}
\newcommand{\stir}[2]{\genfrac{[}{]}{0pt}{}{#1}{#2}}
\newcommand{\src}[1]{\nolinkurl{#1}}
\newcommand{\pin}{\src{b7667ae40a2957eed9a53d8559d53a3ee20dd1c1}}
\newcommand{\zjet}[2]{\zeta^{[#1]}(#2)}
\hypersetup{pdftitle={Continuous-Depth Lerch Calculus},pdfauthor={Research continuation prepared for Vladimir Reshetnikov},pdfsubject={Gamma kernels, exact harmonic identities, Stieltjes convolutions and anchored primitives}}
\title{\textbf{Continuous-Depth Lerch Calculus}\\[5pt]
\large Gamma kernels, all-order harmonic identities,\\
Stieltjes convolutions, and anchored primitives}
\author{Research continuation prepared for Vladimir Reshetnikov's ProveIt programme\\
\small Mathematical derivations and computational checks with OpenAI assistance}
\date{10 October 2026}
\begin{document}
\maketitle
\begin{abstract}
The reciprocal-Lerch continuation asks for a noninteger-depth interpretation
of the multiplier $(\pi\csc\pi p)^\rho(p+a-1)^{-s}$. We answer this question
on $\Rea\rho>0$ by an explicit Gamma-product kernel and a jointly holomorphic
calculus that is entire in the spectral order. Convolution adds both depth
and spectral order. Consequently, an integer \emph{total} depth retains
finite Hurwitz--Lerch closure even when none of the individual depths is
an integer. A polynomial recurrence evaluates every logarithmic moment,
and a local-monodromy argument specifies exactly one finite integer-depth
representation class that cannot contain a noninteger-depth kernel.

The centered kernel has a Fourier--Mellin reflection formula involving a
Pochhammer-weighted alternating Hurwitz series. Its odd values yield an
all-order Gamma-quotient generating identity for absolutely convergent
alternating harmonic sums. Every coefficient is a finite homogeneous
polynomial in ordinary zeta values. Explicit quadratic and cubic harmonic
identities, and a continuous Gamma-weighted odd-polylogarithm family, are
proved. Independent spectral and depth differentiation gives ordinary
Stieltjes-product identities and polygamma-weighted Gamma convolutions;
all shift-primitive orders are obtained with removable resonances explicit.

The Gamma transform and its normalized Meixner semigroup are classical,
not discoveries claimed here. The contribution is their rigorous connection
to the stated ProveIt question and the resulting identity calculus. The
article contains proofs, reproducible finite algebra checks, floating-point
diagnostics, an audit of the inspected source scope, and further research
questions. It does not settle the separate Gaussian $S_6$ or $S_8$ candidates.
\end{abstract}
\vspace{5pt}
\noindent\textbf{Source anchor.} Current repository reading is pinned to \pin.
The earlier \emph{Reciprocal Lerch Convolutions} report was read in its
accessible Library TEX copy, including its precise noninteger-depth
question. No repository or persistent Library files were modified.
\clearpage
\begingroup
\small
\setlength{\parskip}{0pt}
\tableofcontents
\endgroup
\clearpage

\section{Research target, conventions, and contribution}
\label{cdl:sec:scope}

The canonical ProveIt manuscript already develops an extensive system of
polylogarithmic identities, harmonic sums, Hurwitz order derivatives,
Stieltjes distributions, and Gamma correlations \cite{Repo}. A preceding
report, \emph{Reciprocal Lerch Convolutions} \cite{RLC}, proves entire-order
addition at every \emph{integer} multiplicative depth. Its Section 12,
question 10, asks for a real kernel and exact special-function expression
at noninteger depth, together with a precisely delimited explanation of
integer-depth finite closure. This is the problem addressed here.

The answer is not obtained by differentiating an integer-index formula.
We first construct a holomorphic depth parameter from the classical
Gamma transform. Only then do we differentiate, convolve, and specialize.
This order of operations is essential, especially at negative spectral
orders and at collisions with zeta poles.

\subsection{The principal formulas}
For $\Rea\rho>0$ and real $x$, define
\begin{equation}\label{cdl:eq:K}
 K_\rho(x)=\frac{(2\pi)^{\rho-2}e^{x/2}}{\Gamma(\rho)}
 \Gamma\left(\frac\rho2+\frac{ix}{2\pi}\right)
 \Gamma\left(\frac\rho2-\frac{ix}{2\pi}\right).
\end{equation}
If $\rho$ is real and positive, the Gamma product is a squared modulus.
For complex $\rho$ it is a product, not a squared modulus. For
$\Rea a>0$, its spectral extension is
\begin{equation}\label{cdl:eq:Fintro}
 \F_{\rho,a}(s;x)=\frac1{2\pi i}
 \int_{c-i\infty}^{c+i\infty}e^{px}
 \left(\frac\pi{\sin\pi p}\right)^\rho(p+a-1)^{-s}\dd p,
 \qquad \max(0,1-\Rea a)<c<1.
\end{equation}
The central addition law is
\begin{equation}\label{cdl:eq:addintro}
 \F_{\rho,a}(s;\cdot)*\F_{\sigma,a}(t;\cdot)
 =\F_{\rho+\sigma,a}(s+t;\cdot),
 \qquad \Rea\rho,\Rea\sigma,\Rea a>0,
\end{equation}
with ordinary absolutely convergent convolutions at real arguments.
At depth one,
\begin{equation}\label{cdl:eq:Lerchintro}
 \F_{1,a}(s;x)=e^x\Phi(-e^x,s,a),\qquad
 \F_{1,1}(s;x)=-\Li_s(-e^x).
\end{equation}
At total depth two the earlier central formula remains
\begin{equation}\label{cdl:eq:E}
 \F_{2,a}(s;0)=\E_a(s),\qquad
 \E_a(s)=s\zeta(s+1,a),\qquad \E_a(0)=1.
\end{equation}
Thus \eqref{cdl:eq:addintro} supplies continuous-depth factorizations of
an entire Hurwitz family that previously arose from integer-depth factors.

A separate outcome, proved in Section~\ref{cdl:sec:harmonic}, is the
coefficient identity
\begin{equation}\label{cdl:eq:harmonicintro}
 \begin{split}
 &\sum_{k=1}^\infty\frac{(-1)^{k-1}}{k^{2m}}
 \prod_{j=1}^k\left(1+\frac\epsilon j\right)
 \left(1+\frac\epsilon{2k}\right)^{-2m-1}\\
 &\hspace{15mm}=\frac12[z^{2m}]
 \frac{\Gamma(1+\epsilon/2+z)\Gamma(1+\epsilon/2-z)}
      {\Gamma(1+\epsilon)},\qquad m\ge1,\quad |\epsilon|<\tfrac12.
 \end{split}
\end{equation}
Every derivative at $\epsilon=0$ is an absolutely convergent harmonic
identity; the right side reduces by a terminating coefficient calculation.

\subsection{Notation and analytic branches}
The ordinary Lerch series is
$\Phi(z,s,a)=\sum_{n\ge0}z^n(n+a)^{-s}$ in its initial domain.
We use its usual continuation and the right-half-plane logarithm of
$n+a$ when $\Rea a>0$ \cite{DLMFLerch,DLMFZeta}.
The generalized Stieltjes constants are normalized by
\begin{equation}\label{cdl:eq:stieltjes}
 \zeta(1+w,a)=\frac1w+
 \sum_{j=0}^\infty\frac{(-1)^j\gamma_j(a)}{j!}w^j,
 \qquad \gamma_0(a)=-\psi(a).
\end{equation}
We write $\zeta^{[j]}(s,a)=\partial_s^j\zeta(s,a)$ to distinguish
spectral differentiation from differentiation in $a$. The rising factorial
is $(s)_n=s(s+1)\cdots(s+n-1)$, with $(s)_0=1$.
The complete exponential Bell polynomials are fixed by
\begin{equation}\label{cdl:eq:Bell}
 \exp\left(\sum_{j\ge1}d_j\frac{u^j}{j!}\right)
 =\sum_{q\ge0}\B_q(d_1,\ldots,d_q)\frac{u^q}{q!}.
\end{equation}
For example $\B_2=d_1^2+d_2$ and
$\B_3=d_1^3+3d_1d_2+d_3$.

Throughout $g(p)=\pi/\sin\pi p$. On $0<\Rea p<1$ its real part is
positive, because the real part of $\sin\pi p$ is positive. Consequently
$g(p)^\rho=\exp(\rho\log g(p))$ has an unambiguous holomorphic logarithm,
chosen real on $0<p<1$. The logarithm of $p+a-1$ in
\eqref{cdl:eq:Fintro} is the right-half-plane logarithm.

\subsection{What is classical, and what is established here}
The Fourier transform behind \eqref{cdl:eq:K} is Ramanujan's classical
Gamma-product formula; a modern proof and complex-parameter statement
are given by Chakrabarti and Srinivasan \cite{CS}. After centering and
normalization, $K_\rho$ is the classical symmetric Meixner semigroup
\cite{Schoutens}. The elementary beta and Gamma identities used below
are also classical \cite{DLMFBeta,DLMFGamma}.

We claim a rigorous continuation of the specified source question, not
historical priority for these inputs. The main additions relative to the
inspected reciprocal report are the continuous-depth/spectral family,
its complete finite moment calculus, continuous-split Stieltjes identities,
the harmonic coefficient master, and the explicit representation-class
obstruction. Some low-weight harmonic specializations recover familiar
Euler sums. No assertion that every printed special case is previously
unpublished is required by, or inferred from, the proofs.

\section{The Gamma kernel and its holomorphic semigroup}
\label{cdl:sec:kernel}

\begin{lemma}[Beta Fourier identity]\label{cdl:lem:beta}
For $\Rea\rho>0$ and real $v$,
\begin{equation}\label{cdl:eq:betaFT}
 \frac{\Gamma(\rho/2+iv)\Gamma(\rho/2-iv)}{\Gamma(\rho)}
 =\int_\R e^{ivy}\bigl(2\cosh(y/2)\bigr)^{-\rho}\dd y.
\end{equation}
Both sides are holomorphic in $\rho$ on this half-plane.
\end{lemma}
\begin{proof}
Euler's beta integral on $(0,\infty)$ gives the left side as
$\int_0^\infty t^{\rho/2+iv-1}(1+t)^{-\rho}\dd t$.
Set $t=e^y$ and use
$e^{\rho y/2}(1+e^y)^{-\rho}=(2\cosh(y/2))^{-\rho}$.
On compact subsets of $\Rea\rho>0$, any fixed parameter derivative is
bounded by a polynomial in $|y|$ times an exponentially decreasing
function. This proves convergence and holomorphy.
\end{proof}

\begin{theorem}[Continuous-depth transform]\label{cdl:thm:kernel}
For $\Rea\rho>0$ and $0<\Rea p<1$,
\begin{equation}\label{cdl:eq:Ktransform}
 \int_\R e^{-px}K_\rho(x)\dd x=g(p)^\rho.
\end{equation}
For every real $c\in(0,1)$ the function $e^{-cx}K_\rho(x)$ belongs to the
Schwartz space and depends holomorphically on $\rho$ in its Schwartz
seminorms. In particular,
\begin{equation}\label{cdl:eq:Ksemi}
 \int_\R K_\rho(y)K_\sigma(x-y)\dd y=K_{\rho+\sigma}(x)
\end{equation}
is an ordinary absolutely convergent identity for real $x$ and positive
real parts of $\rho,\sigma$.
\end{theorem}
\begin{proof}
Fourier inversion in \eqref{cdl:eq:betaFT}, with the change of scale
$x=2\pi v$, gives
\[
 e^{-x/2}K_\rho(x)=\frac1{2\pi}
 \int_\R e^{itx}\left(\frac\pi{\cosh\pi t}\right)^\rho\dd t.
\]
The Gamma product and the hyperbolic kernel are integrable and smooth;
Gamma asymptotics justify the ordinary inversion. This is precisely
\eqref{cdl:eq:Ktransform} on $\Rea p=1/2$. The same asymptotics show
that $K_\rho(x)$ is bounded by a fixed polynomial in $1+|x|$ as
$x\to+\infty$, and by such a polynomial times $e^x$ as
$x\to-\infty$, locally uniformly in $\rho$. Argument and depth
derivatives obey analogous bounds, with possible logarithmic factors.
Multiplication by $e^{-cx}$ therefore gives all Schwartz estimates.
Holomorphy and the identity theorem extend the transform throughout
$0<\Rea p<1$.

For the convolution put $k_\rho(y)=e^{-cy}K_\rho(y)$. Its absolute
integral is bounded by
$e^{cx}\|k_\rho\|_1\|k_\sigma\|_\infty$.
The product of their Fourier transforms is
$g(c+it)^{\rho+\sigma}$, because the same logarithm was used in both
factors. Fourier inversion proves \eqref{cdl:eq:Ksemi} pointwise.
\end{proof}

The endpoint normalizations are
\begin{equation}\label{cdl:eq:seeds}
 K_1(x)=\frac1{1+e^{-x}},\qquad
 K_2(x)=\frac{x}{1-e^{-x}},\qquad K_2(0)=1.
\end{equation}
They follow from Gamma reflection and recurrence. Define
$b_\rho(x)=e^{-x/2}K_\rho(x)$. For real $\rho>0$, the density
$b_\rho(x)/\pi^\rho$ is positive, has integral one, and has characteristic
function $\operatorname{sech}(\pi t)^\rho$. In Schoutens's notation it is
$\operatorname{Meixner}(2\pi,0,\rho/2,0)$ \cite{Schoutens}.
This identifies a classical semigroup; the weighted uncentered convention
is what makes it fit the Lerch profiles.

\begin{remark}[Complex argument strip]
The Gamma formula is holomorphic for
$|\Ima x|<\pi\Rea\rho$. Its first boundary poles occur at
$x=\pm i\pi\rho$, and further poles at $x=\pm i\pi(\rho+2n)$.
Thus this symmetric horizontal strip cannot be enlarged for the base
kernel. The pole description concerns the argument $x$, whereas the
monodromy obstruction of Section~\ref{cdl:sec:obstruction} concerns the
transform variable $p$; these are different assertions.
\end{remark}

\section{Entire spectral orders and two-parameter addition}
\label{cdl:sec:spectral}

\begin{theorem}[Entire spectral family]\label{cdl:thm:spectral}
Let $\Rea\rho>0$ and $\Rea a>0$. Formula
\eqref{cdl:eq:Fintro} defines an entire function of $s\in\C$, independent
of the permitted choice of $c$. It is locally jointly holomorphic in
$\rho,a,s$, and holomorphic in $x$ on $|\Ima x|<\pi\Rea\rho$.
For real $x$ and every permitted $c$, the map
$(\rho,a,s)\mapsto e^{-cx}\F_{\rho,a}(s;x)$ is locally holomorphic into
Schwartz space. All fixed polynomial moments and all fixed mixed
parameter derivatives are therefore integrable and bounded.

If $\Rea s>0$, then
\begin{equation}\label{cdl:eq:fractional}
 \F_{\rho,a}(s;x)=\frac1{\Gamma(s)}\int_0^\infty
 u^{s-1}e^{-(a-1)u}K_\rho(x-u)\dd u.
\end{equation}
On the entire spectral plane,
\begin{equation}\label{cdl:eq:ladder}
 (\partial_x+a-1)\F_{\rho,a}(s;x)=\F_{\rho,a}(s-1;x).
\end{equation}
\end{theorem}
\begin{proof}
On a vertical line in the permitted strip, $g(c+it)^\rho$ decays as
$e^{-\pi\Rea\rho|t|}$ times a locally bounded factor. The argument of
$g$ is bounded, and $(c+it+a-1)^{-s}$ grows at most polynomially on
compact parameter sets. Derivatives in $s$ insert powers of
$-\log(c+it+a-1)$; derivatives in $\rho$ insert powers of
$\log g(c+it)$, which have at most linear growth. Derivatives in $a$
and $t$ preserve exponential decay up to polynomial factors. Consequently
the multiplier is a holomorphic Schwartz family. Fourier inversion,
continuous on Schwartz space, proves the norm assertions, as well as
pointwise holomorphy. The factor $e^{itx}$ gives the stated argument strip.

Cauchy's theorem on rectangles between two allowed vertical lines proves
independence of $c$: the horizontal integrals tend to zero by the same
exponential estimate, and no pole or branch boundary lies between them.
For $\Rea s>0$, the Fourier transform of the weighted right side of
\eqref{cdl:eq:fractional} is the product of \eqref{cdl:eq:Ktransform}
and the elementary Laplace integral $(p+a-1)^{-s}$. Fubini is justified
by $c+\Rea a-1>0$, the weighted kernel estimate, and the integrability
of $u^{\Rea s-1}$ at zero. Transform uniqueness gives the identity.
Finally, multiplying the contour integrand by $p+a-1$ proves
\eqref{cdl:eq:ladder} for all $s$ directly.
\end{proof}

\begin{corollary}[Lerch recovery and negative orders]\label{cdl:cor:recovery}
The identifications \eqref{cdl:eq:Lerchintro} hold for all complex $s$.
Moreover, for every nonnegative integer $N$,
\begin{equation}\label{cdl:eq:negative}
 \F_{\rho,a}(0;x)=K_\rho(x),\qquad
 \F_{\rho,a}(-N;x)=(\partial_x+a-1)^N K_\rho(x).
\end{equation}
\end{corollary}
\begin{proof}
At $\rho=1$, insert the first seed in \eqref{cdl:eq:fractional} and
compare with the Euler integral for $e^x\Phi(-e^x,s,a)$.
Both functions are entire in $s$, so equality extends from $\Rea s>0$.
The polylogarithm specialization is the defining Lerch relation.
The negative-order formulas follow from the contour at $s=0$ and the ladder.
\end{proof}

In particular, one must not insert $s=-N$ into the improper integral
\eqref{cdl:eq:fractional} and set $1/\Gamma(-N)=0$.
The integral representation has left its convergence domain; the true
value is the generally nonzero differential expression
\eqref{cdl:eq:negative}.

\begin{theorem}[Simultaneous addition and differentiated products]
\label{cdl:thm:addition}
Let $\Rea a>0$, $\Rea\rho_j>0$, $s_j\in\C$, and $x\in\R$.
For any finite number $r\ge2$ of factors,
\begin{equation}\label{cdl:eq:addition}
 \int_{\R^{r-1}}\prod_{j=1}^r\F_{\rho_j,a}(s_j;x_j)
 \dd x_1\cdots\dd x_{r-1}
 =\F_{R,a}(W;x),
 \quad R=\sum_j\rho_j,\quad W=\sum_js_j,
\end{equation}
where $x_r=x-\sum_{j<r}x_j$. The integral converges absolutely.
Every fixed mixed depth derivative and spectral derivative passes through
it. A derivative of total depth degree $q$ and total spectral degree $M$
is $\partial_R^q\partial_W^M\F_{R,a}(W;x)$.
\end{theorem}
\begin{proof}
Choose one common permitted $c$. The product of the exponential weights
on the hyperplane is $e^{-cx}$, a constant. The absolute integral is at
most $e^{cx}$ times the product of $r-1$ weighted $L^1$ norms and one
weighted $L^\infty$ norm. Theorem~\ref{cdl:thm:spectral} gives locally
uniform versions of these estimates after all the stated derivatives.
Multiplication of the Fourier transforms yields
$g(p)^R(p+a-1)^{-W}$. Fourier inversion proves the identity.
The norm holomorphy proves the derivative statements, not merely formal
differentiation of the multipliers.
\end{proof}

The number of factors remains an integer, but their \emph{depth
parameters} need not. This distinction is the resolution of the source
question. No fractional-dimensional integration domain is introduced.
The entire function on the right continues the result to its full
argument strip, even when a particular complex-argument integral on the
left would have a smaller common strip.

\section{Finite depth shifts and generalized harmonic coefficients}
\label{cdl:sec:shifts}

\begin{theorem}[A continuous two-step recurrence]\label{cdl:thm:shift}
For $\Rea\rho>0$,
\begin{equation}\label{cdl:eq:Kshift}
 \rho(\rho+1)K_{\rho+2}(x)=(x^2+\pi^2\rho^2)K_\rho(x).
\end{equation}
For $\Rea a>0$ and every complex $s$,
\begin{equation}\label{cdl:eq:Fshift}
 \begin{split}
 \rho(\rho+1)\F_{\rho+2,a}(s;x)
 ={}&(x^2+\pi^2\rho^2)\F_{\rho,a}(s;x)\\
 &-2xs\F_{\rho,a}(s+1;x)
   +s(s+1)\F_{\rho,a}(s+2;x).
 \end{split}
\end{equation}
More generally, put
\begin{equation}\label{cdl:eq:shiftpoly}
 Q_n(\rho,X)=\frac{\prod_{j=0}^{n-1}(X^2+\pi^2(\rho+2j)^2)}{(\rho)_{2n}}.
\end{equation}
Then
\begin{equation}\label{cdl:eq:allshifts}
 \F_{\rho+2n,a}(s;x)=
 \sum_{k=0}^{2n}\frac{(-1)^k\partial_X^kQ_n(\rho,x)}{k!}
 (s)_k\F_{\rho,a}(s+k;x).
\end{equation}
\end{theorem}
\begin{proof}
Gamma recurrence in the two numerator factors of \eqref{cdl:eq:K}
multiplies their product by
$(\rho/2+ix/(2\pi))(\rho/2-ix/(2\pi))$.
The prefactor changes by $4\pi^2/[\rho(\rho+1)]$, proving
\eqref{cdl:eq:Kshift}. Iteration gives
$K_{\rho+2n}(x)=Q_n(\rho,x)K_\rho(x)$.
For $\Rea s>0$, substitute this identity at $x-u$ in
\eqref{cdl:eq:fractional}. Expand the polynomial by its finite Taylor
formula. Its $u^k$ term changes $1/\Gamma(s)$ to
$\Gamma(s+k)/\Gamma(s)=(s)_k$ in the remaining integral.
This proves \eqref{cdl:eq:allshifts}, with the displayed two-step case.
Entire continuation in $s$ finishes the proof. Apparent zeros of a
Pochhammer factor are interpreted in the complete entire expression.
\end{proof}

The finite coefficient family has a useful harmonic description.
For $N\ge0$ and away from denominator poles, write
\[
 H_N^{(j)}(z)=\sum_{\ell=0}^{N-1}(z+\ell)^{-j}.
\]
The factorization
\begin{equation}\label{cdl:eq:PochhammerQ}
 Q_n(\rho,X)=\frac{(4\pi^2)^n
 (\rho/2+iX/(2\pi))_n(\rho/2-iX/(2\pi))_n}{(\rho)_{2n}}
\end{equation}
yields
\begin{equation}\label{cdl:eq:Qjets}
 \begin{split}
 \partial_\rho^j\log Q_n(\rho,X)
 ={}&(-1)^{j-1}(j-1)!\biggl[
 2^{-j}H_n^{(j)}\left(\frac\rho2+\frac{iX}{2\pi}\right)\\
 &\quad+2^{-j}H_n^{(j)}\left(\frac\rho2-\frac{iX}{2\pi}\right)
       -H_{2n}^{(j)}(\rho)\biggr].
 \end{split}
\end{equation}
Bell polynomials therefore express every depth derivative of $Q_n$ as
finite generalized-harmonic combinations. If parameters are continued
to a zero of $Q_n$, one differentiates the product
\eqref{cdl:eq:shiftpoly} rather than dividing by it; the logarithmic
formula is only an auxiliary representation away from such zeros.

\subsection{Recovery of the integer finite basis}
For integers $m\ge1$, define
\begin{align}
 p_{2m}(X)&=X\prod_{j=1}^{m-1}(X^2+4j^2\pi^2),\label{cdl:eq:pEven}\\
 p_{2m+1}(X)&=\prod_{j=1}^{m}(X^2+(2j-1)^2\pi^2).
 \label{cdl:eq:pOdd}
\end{align}
The seeds \eqref{cdl:eq:seeds} and the recurrence imply the existing
integer-depth formula
\begin{equation}\label{cdl:eq:integerK}
 K_r(x)=\frac{p_r(x)}{(r-1)!\{1-(-1)^r e^{-x}\}}.
\end{equation}
At even depth and $x=0$, take the continuous value.
This is the central-factorial family in \cite{RLC}; its recovery is a
normalization check rather than a new claim.

If
$p_{2m}(X)=\sum_k A_{m,k}X^{2k+1}$ and
$p_{2m+1}(X)=\sum_k B_{m,k}X^{2k}$, then
\begin{align}
 \F_{2m,a}(s;0)&=\frac1{(2m-1)!}\sum_{k=0}^{m-1}
 A_{m,k}(s)_{2k}\E_a(s+2k),\label{cdl:eq:evenbasis}\\
 \F_{2m+1,a}(s;0)&=\frac1{(2m)!}\sum_{k=0}^{m}
 B_{m,k}(s)_{2k}\eta(s+2k,a),\label{cdl:eq:oddbasis}
\end{align}
where
\begin{equation}\label{cdl:eq:eta}
 \eta(s,a)=\Phi(-1,s,a)
 =2^{-s}\left\{\zeta(s,a/2)-\zeta(s,(a+1)/2)\right\}
\end{equation}
is entire in $s$. To verify these formulas directly, insert
\eqref{cdl:eq:integerK} in the positive-order integral at zero and
expand $(e^u\mp1)^{-1}$. In even depth the factor
$(s)_{2k+1}\zeta(s+2k+1,a)$ is best written
$(s)_{2k}\E_a(s+2k)$; this displays its removable pole before
continuation. In particular,
\begin{equation}\label{cdl:eq:C4}
 \F_{4,a}(s;0)=\frac{s(s+1)\E_a(s+2)+4\pi^2\E_a(s)}6.
\end{equation}
The same reasoning away from zero gives finite ordinary Lerch values.
For $x<0$, an explicit form is
\begin{equation}\label{cdl:eq:finiteLerch}
 \F_{r,a}(s;x)=\frac{-\varepsilon_r e^x}{(r-1)!}
 \sum_{j=0}^{r-1}\frac{(-1)^jp_r^{(j)}(x)}{j!}
 (s)_j\Phi(\varepsilon_r e^x,s+j,a),\qquad
 \varepsilon_r=(-1)^r.
\end{equation}
Other real arguments are reached by continuing the \emph{whole}
combination. At even depth the individual Lerch terms can cross a cut;
this does not authorize independent changes of their branches.

\section{Centered reflection and exact Gamma--polylogarithm values}
\label{cdl:sec:reflection}

\subsection{A Pochhammer-weighted alternating Hurwitz family}
For $\Rea\rho>0$, initially $\Rea s>\Rea\rho$, define
\begin{equation}\label{cdl:eq:Zseries}
 Z_\rho(s)=\sum_{n=0}^\infty
 \frac{(-1)^n(\rho)_n}{n!(n+\rho/2)^s}.
\end{equation}
This is a one-parameter Pochhammer-weighted Hurwitz--Lerch series, not
an ordinary Lerch function for a general $\rho$. We define it explicitly
to avoid dependence on competing generalized-Lerch notation.

\begin{proposition}[Entire continuation and negative values]
\label{cdl:prop:Z}
The function $Z_\rho$ extends to an entire function of $s$. On
$\Rea s>0$ it has the ordinary integral
\begin{equation}\label{cdl:eq:ZMellin}
 Z_\rho(s)=\frac1{\Gamma(s)}\int_0^\infty
 t^{s-1}\bigl(2\cosh(t/2)\bigr)^{-\rho}\dd t.
\end{equation}
Writing $f_\rho(t)=(2\cosh(t/2))^{-\rho}$, one has
\begin{equation}\label{cdl:eq:Znegative}
 Z_\rho(-N)=(-1)^N f_\rho^{(N)}(0),\qquad N=0,1,2,\ldots.
\end{equation}
In particular $Z_\rho(0)=2^{-\rho}$ and
$Z_\rho(-2j-1)=0$ for every $j\ge0$.
\end{proposition}
\begin{proof}
Expand $e^{-\rho t/2}(1+e^{-t})^{-\rho}$ by its binomial series.
In the absolute-convergence region $\Rea s>\Rea\rho$, Fubini follows
from $(\rho)_n/n!=O(n^{\Rea\rho-1})$ on compact parameter sets and
from the Laplace integral with $n+\rho/2$ in the right half-plane.
This proves agreement of \eqref{cdl:eq:Zseries} and
\eqref{cdl:eq:ZMellin}. The latter converges on $\Rea s>0$.

For any chosen $N$, subtract the Taylor polynomial of $f_\rho(t)$ at
zero inside the integral over $(0,1)$ and integrate the polynomial
terms separately. The resulting Mellin transform is meromorphic to
$\Rea s>-N-1$, with simple poles at $s=-j$ of residue
$f_\rho^{(j)}(0)/j!$. Multiplication by $1/\Gamma(s)$ cancels all these
poles, since its zero at $-j$ has derivative $(-1)^j j!$.
This gives \eqref{cdl:eq:Znegative} and entire continuation. The kernel
is even at zero, so all its odd derivatives vanish. The same subtracted
integrals, locally uniform in $\rho$, establish joint holomorphy.
\end{proof}

For later comparison,
\begin{equation}\label{cdl:eq:Zspecial}
 Z_1(s)=2^s\beta(s),\qquad Z_2(s)=\eta(s-1),
\end{equation}
where $\beta$ is Dirichlet's beta function and
$\eta(s)=(1-2^{1-s})\zeta(s)$ is the entire eta function.
These equalities follow from the series in an absolute domain and then
from analytic continuation.

\begin{theorem}[Centered Fourier--Mellin reflection]
\label{cdl:thm:reflection}
For $\Rea\rho>0$ and every complex $w$,
\begin{equation}\label{cdl:eq:reflection}
 \F_{\rho,1/2}(w;0)=
 2(2\pi)^{\rho+w-2}\Gamma(1-w)
 \cos\left(\frac{\pi w}2\right)Z_\rho(1-w),
\end{equation}
where all apparent singularities of the right side are removable.
Equivalently, as an identity of meromorphic functions,
\begin{equation}\label{cdl:eq:completed}
 \Gamma(w)\sin\left(\frac{\pi w}2\right)\F_{\rho,1/2}(w;0)
 =\pi(2\pi)^{\rho+w-2}Z_\rho(1-w).
\end{equation}
\end{theorem}
\begin{proof}
First suppose $0<\Rea w<1$. At $a=1/2$, the fractional representation is
\[\F_{\rho,1/2}(w;0)=\frac1{\Gamma(w)}\int_0^\infty x^{w-1}b_\rho(x)\dd x.\]
The Fourier convention is $\widehat f(v)=\int_\R e^{ivt}f(t)\dd t$.
By \eqref{cdl:eq:betaFT},
\[b_\rho(x)=(2\pi)^{\rho-2}\widehat f_\rho(x/(2\pi)).\]
We need the Mellin transform of this Fourier transform.

Insert a factor $e^{-\delta v}$, $\delta>0$, in
$\int_0^\infty v^{w-1}\widehat f_\rho(v)\dd v$.
Fubini is then absolute. The inner cosine integral is
\[
 \int_0^\infty v^{w-1}e^{-\delta v}\cos(vt)\dd v
 =\frac{\Gamma(w)}2\bigl((\delta-it)^{-w}+(\delta+it)^{-w}\bigr).
\]
It is bounded in modulus by a constant times $t^{-\Rea w}$ for $t>0$,
uniformly as $\delta\downarrow0$. This is integrable against
$f_\rho(t)$ in our strip. Dominated convergence therefore gives
\[
 \int_0^\infty v^{w-1}\widehat f_\rho(v)\dd v
 =2\Gamma(w)\cos(\pi w/2)
   \int_0^\infty t^{-w}f_\rho(t)\dd t.
\]
On the left the regulator may be removed because the Gamma product
has exponential decay. Rescaling and using \eqref{cdl:eq:ZMellin}
proves \eqref{cdl:eq:reflection}. Both sides have meromorphic
continuations in $w$, while the left side is entire; uniqueness proves
the global identity and removability. Euler reflection for Gamma gives
\eqref{cdl:eq:completed}.
\end{proof}

The removability can be seen concretely. For $m\ge1$,
\begin{equation}\label{cdl:eq:evenreflection}
 \F_{\rho,1/2}(2m;0)=
 \frac{2(-1)^{m+1}(2\pi)^{\rho+2m-2}}{(2m-1)!}
 Z_\rho^{[1]}(1-2m).
\end{equation}
At $w=2m+1$, $m\ge0$, the zero of the cosine cancels the Gamma pole,
giving
\begin{equation}\label{cdl:eq:oddreflection}
 \F_{\rho,1/2}(2m+1;0)=
 \frac{(-1)^m\pi(2\pi)^{\rho+2m-1}}{(2m)!}Z_\rho(-2m).
\end{equation}
For example, at $\rho=2,w=2$ this becomes the standard relation
between $\zeta'(-2)$ and $\zeta(3)$, with the half-shift factor intact.
These formulas also explain why substituting a trivial zero before
combining a pole gives a wrong answer.

\subsection{Two odd-order evaluation families}
\begin{theorem}[Odd values and Gamma coefficients]\label{cdl:thm:oddZ}
For every $m\ge0$ and $\Rea\rho>0$,
\begin{align}
 Z_\rho(2m+1)
 &=\frac1{2\Gamma(\rho)}[z^{2m}]
   \Gamma(\rho/2+z)\Gamma(\rho/2-z)\label{cdl:eq:oddZ}\\
 &=\frac{\Gamma(\rho/2)^2}{2\Gamma(\rho)}[z^{2m}]
 \exp\left(\sum_{j\ge1}\frac{\zeta(2j,\rho/2)}jz^{2j}\right).
 \label{cdl:eq:oddZexp}
\end{align}
Thus all these values lie in the Gamma-prefactor times the polynomial
ring generated by the finitely many Hurwitz values
$\zeta(2,\rho/2),\ldots,\zeta(2m,\rho/2)$.
\end{theorem}
\begin{proof}
The beta integral gives, for $|\Rea z|<\Rea\rho/2$,
\[
 \int_0^\infty\cosh(zt)f_\rho(t)\dd t
 =\frac{\Gamma(\rho/2+z)\Gamma(\rho/2-z)}{2\Gamma(\rho)}.
\]
The coefficient of $z^{2m}$ on the left is precisely
\eqref{cdl:eq:ZMellin} at $s=2m+1$. The logarithmic Taylor expansion
of the Gamma product, using
$\psi^{(2j-1)}(h)=(2j-1)!\zeta(2j,h)$, gives
\eqref{cdl:eq:oddZexp}.
\end{proof}

The first three values are the prefactor in \eqref{cdl:eq:oddZexp}
times $1$, $\zeta(2,\rho/2)$, and
$\tfrac12(\zeta(2,\rho/2)^2+\zeta(4,\rho/2))$.
Where the original series \eqref{cdl:eq:Zseries} is not convergent,
these are values of its proved continuation, not assertions about a
divergent ordinary sum.

Define the finite polynomials
\begin{equation}\label{cdl:eq:Emompoly}
 U_m(\rho)=[z^{2m}]\sec(z)^\rho.
\end{equation}
They are polynomials in $\rho$ with rational coefficients. For example,
\begin{equation}\label{cdl:eq:Umexamples}
 U_0=1,\quad U_1=\frac\rho2,\quad
 U_2=\frac{\rho(3\rho+2)}{24},\quad
 U_3=\frac{\rho(15\rho^2+30\rho+16)}{720}.
\end{equation}
The negative-value formula in Proposition~\ref{cdl:prop:Z}, inserted
in \eqref{cdl:eq:oddreflection}, proves
\begin{equation}\label{cdl:eq:oddF}
 \boxed{\displaystyle
 \F_{\rho,1/2}(2m+1;0)=\frac{\pi^{\rho+2m}}2U_m(\rho).}
\end{equation}
Indeed,
$f_\rho(t)=2^{-\rho}\operatorname{sech}(t/2)^\rho$, whose $t^{2m}$
coefficient is $2^{-\rho-2m}(-1)^mU_m(\rho)$.
This second family is polynomial in depth, apart from its explicit
power of $\pi$; it is distinct from the Hurwitz-polygamma coefficients
in \eqref{cdl:eq:oddZexp}.

\begin{corollary}[A continuous Gamma-weighted polylogarithm identity]
\label{cdl:cor:polylog}
Let $\Rea\rho>0$ and $m\ge0$. Put
$G_\rho(x)=\Gamma(\rho/2+ix/(2\pi))\Gamma(\rho/2-ix/(2\pi))$.
Then the ordinary integral
\begin{equation}\label{cdl:eq:polylogGamma}
 \begin{split}
 &\int_\R G_\rho(x)
 \frac{\Li_{2m+1}(ie^{-x/2})-\Li_{2m+1}(-ie^{-x/2})}{i}\dd x\\
 &\hspace{15mm}=
 \frac{\Gamma(\rho)\pi^{2m+3}}{2^{\rho+2m-1}}U_m(\rho+1)
 \end{split}
\end{equation}
is absolutely convergent, with the usual principal polylogarithms.
\end{corollary}
\begin{proof}
Separating the odd terms of the polylogarithm series gives
\[
 \F_{1,1/2}(s;x)=\frac{2^{s-1}e^{x/2}}i
 \{\Li_s(ie^{x/2})-\Li_s(-ie^{x/2})\}.
\]
It holds first for $x<0$, then along the real axis by continuation;
no argument crosses the positive polylogarithm cut. Convolve this
profile with $K_\rho$ at zero. Theorem~\ref{cdl:thm:addition}
gives $\F_{\rho+1,1/2}(2m+1;0)$, evaluated by
\eqref{cdl:eq:oddF}. Substituting the Gamma prefactor and collecting
powers of $2$ and $\pi$ yields \eqref{cdl:eq:polylogGamma}.
Absolute convergence is already part of the convolution theorem.
\end{proof}

For $m=0$ the polylogarithm odd part is
$2\arctan(e^{-x/2})$. This base case also follows by the evenness of
$G_\rho$ and the identity
$\arctan(e^u)+\arctan(e^{-u})=\pi/2$.
For $m=1$ the right side of \eqref{cdl:eq:polylogGamma} is
$\Gamma(\rho)(\rho+1)\pi^5/2^{\rho+2}$.
Differentiation in $\rho$, justified by the kernel estimates, generates
polygamma-weighted versions of every member of this family.

\section{An all-order alternating harmonic master}
\label{cdl:sec:harmonic}

Write $H_k^{(j)}=\sum_{\ell=1}^k\ell^{-j}$ and $H_k=H_k^{(1)}$.
The next theorem converts the centered continuous-depth calculation
into an ordinary absolutely convergent Euler-sum family.

\begin{theorem}[Gamma-quotient harmonic generating function]
\label{cdl:thm:harmonicmaster}
For every integer $m\ge1$ and $|\epsilon|<1/2$, identity
\eqref{cdl:eq:harmonicintro} holds. Its series is locally normally
convergent in $\epsilon$, and every derivative can be taken termwise.
If
\begin{equation}\label{cdl:eq:dj}
 d_j(k,m)=(-1)^{j-1}(j-1)!
 \left(H_k^{(j)}-\frac{2m+1}{2^jk^j}\right),
\end{equation}
then, for every $q\ge0$,
\begin{equation}\label{cdl:eq:mastercoeff}
 \begin{split}
 &\sum_{k=1}^\infty\frac{(-1)^{k-1}}{k^{2m}}
 \B_q\bigl(d_1(k,m),\ldots,d_q(k,m)\bigr)\\
 &\quad=\frac{q!}{2}[\epsilon^qz^{2m}]
 \exp\left(\sum_{\ell\ge2}\frac{(-1)^\ell\zeta(\ell)}\ell
 \left\{(\epsilon/2+z)^\ell+(\epsilon/2-z)^\ell-\epsilon^\ell\right\}\right).
 \end{split}
\end{equation}
Only terms with $\ell\le2m+q$ are needed on the right. Its value is
a homogeneous polynomial in zeta values of total weight $2m+q$.
\end{theorem}
\begin{proof}
Set $\rho=2+\epsilon$ in \eqref{cdl:eq:Zseries}, put $k=n+1$, and
multiply by $1+\epsilon$. The elementary equality
\[
 (1+\epsilon)\frac{(2+\epsilon)_{k-1}}{(k-1)!}
 =k\prod_{j=1}^k(1+\epsilon/j)
\]
shows that the left side of \eqref{cdl:eq:harmonicintro} is
$(1+\epsilon)Z_{2+\epsilon}(2m+1)$. Theorem~\ref{cdl:thm:oddZ}
and $\Gamma(2+\epsilon)=(1+\epsilon)\Gamma(1+\epsilon)$ prove
its stated Gamma quotient.

On any compact disk $|\epsilon|\le\delta<1/2$, the finite product is
bounded by $C k^{\Rea\epsilon}$, uniformly, by its Gamma-quotient
expression. The other factor is uniformly bounded. Hence the series
is dominated by $Ck^{-2m+\delta}$, a summable sequence. Derivatives
insert at most fixed powers of $\log k$, so they remain normally
convergent. This also proves absolute convergence of every coefficient
identity without relying on alternation.

The logarithm of the $k$th parameter factor has $j$th derivative at
zero equal to \eqref{cdl:eq:dj}. Its $q$th derivative is therefore
$\B_q(d_1,\ldots,d_q)$ by \eqref{cdl:eq:Bell}.
For the Gamma quotient use the convergent expansion
\[
 \log\Gamma(1+u)=-\gamma u+
 \sum_{\ell\ge2}\frac{(-1)^\ell\zeta(\ell)}\ell u^\ell.
\]
The Euler-constant terms cancel because
$(\epsilon/2+z)+(\epsilon/2-z)-\epsilon=0$.
This proves \eqref{cdl:eq:mastercoeff}. Each displayed term has total
degree $\ell$ in $\epsilon,z$ and zeta weight $\ell$. Exponentiation
and extraction of degree $q+2m$ prove finiteness and homogeneity.
No independence assertion concerning zeta values is used.
\end{proof}

\subsection{Explicit low-order identities}
At $m=1,q=1$ the shifted sum is $-\zeta(3)/2$. After moving the
$k^{-3}$ term, one recovers
\begin{equation}\label{cdl:eq:classicharmonic}
 \sum_{k=1}^\infty\frac{(-1)^{k-1}H_k}{k^2}=\frac58\zeta(3).
\end{equation}
At $m=2,q=1$ the same operation gives
\begin{equation}\label{cdl:eq:weight5}
 \sum_{k=1}^\infty\frac{(-1)^{k-1}H_k}{k^4}
 =\frac{59}{32}\zeta(5)-\frac12\zeta(2)\zeta(3).
\end{equation}
These low-weight recoveries are useful checks, not priority claims.
The next two formulas keep the harmonic combinations naturally produced
by the generating function:
\begin{equation}\label{cdl:eq:harmonicquadratic}
 \boxed{\displaystyle
 \sum_{k=1}^\infty\frac{(-1)^{k-1}}{k^2}
 \left(H_k^2-H_k^{(2)}-\frac{3H_k}{k}+\frac3{k^2}\right)
 =\frac18\zeta(4).}
\end{equation}
\begin{equation}\label{cdl:eq:harmoniccubic}
 \boxed{\begin{aligned}
 &\sum_{k=1}^\infty\frac{(-1)^{k-1}}{k^2}
 \biggl[H_k^3-3H_kH_k^{(2)}+2H_k^{(3)}
 -\frac9{2k}\{H_k^2-H_k^{(2)}\}\\[-2pt]
 &\hspace{42mm}+\frac{9H_k}{k^2}-\frac{15}{2k^3}\biggr]
 =\frac32\{\zeta(2)\zeta(3)-\zeta(5)\}.
 \end{aligned}}
\end{equation}
To make the coefficient calculation transparent, the $z^2$ coefficient
on the right of \eqref{cdl:eq:harmonicintro} is
\begin{equation}\label{cdl:eq:z2master}
 \frac12\frac{\Gamma(1+\epsilon/2)^2}{\Gamma(1+\epsilon)}
 \zeta(2,1+\epsilon/2).
\end{equation}
Here
\begin{align*}
 \log\frac{\Gamma(1+\epsilon/2)^2}{\Gamma(1+\epsilon)}
 &=-\frac{\zeta(2)}4\epsilon^2+
    \frac{\zeta(3)}4\epsilon^3+O(\epsilon^4),\\
 \zeta(2,1+\epsilon/2)
 &=\zeta(2)-\zeta(3)\epsilon+
   \frac34\zeta(4)\epsilon^2-\frac12\zeta(5)\epsilon^3+O(\epsilon^4).
\end{align*}
Multiplying and taking derivatives of order two and three gives
\eqref{cdl:eq:harmonicquadratic} and \eqref{cdl:eq:harmoniccubic};
use $\zeta(2)^2=5\zeta(4)/2$ in the quadratic case.
The left sides follow by expanding $\B_2$ and $\B_3$ with
\eqref{cdl:eq:dj}, fixing every sign and denominator explicitly.

\subsection{A terminating symbolic implementation}
To compute a desired pair $(m,q)$, truncate the exponent in
\eqref{cdl:eq:mastercoeff} at total degree $2m+q$, retaining only
$\epsilon$-degree at most $q$ and $z$-degree at most $2m$.
If $L_{ij}$ is the coefficient of $\epsilon^iz^j$ in that exponent
and $A_{uv}$ the corresponding coefficient of its exponential, then
\begin{equation}\label{cdl:eq:coeffrecurrence}
 A_{00}=1,\qquad
 A_{uv}=\frac1{u+v}\sum_{\substack{i\le u,\ j\le v\\i+j>0}}
 (i+j)L_{ij}A_{u-i,v-j}.
\end{equation}
Thus the answer is $q!A_{q,2m}/2$ using exact rational arithmetic
and formal symbols for the zeta values. The supplied code compares
this bivariate recurrence with a separate computation from
\eqref{cdl:eq:oddZexp}, expanding the Hurwitz values in $\epsilon$.
It generates all $1\le m\le4$, $0\le q\le6$ cases.

This gives a family of exact combinations at arbitrarily large harmonic
degree. It does not claim that an arbitrary individual alternating Euler
sum reduces to ordinary zeta products, nor does it permit differentiation
in the discrete integer $m$. Independent spectral derivatives require
the entire family constructed earlier, not differentiation of a finite
integer-index formula.

\section{Every logarithmic moment by a finite recurrence}
\label{cdl:sec:moments}

The variable $x=\log y$ converts powers of $x$ into logarithmic moments
in the multiplicative variable. The purpose of this section is a finite
algorithm valid at arbitrary complex spectral orders and continuous
depths, not just isolated low moments.

Put $A=p+a-1$, $v=\pi\cot\pi p$, and $h=A^{-1}$.
Define polynomials in $v,h$, with parameters $\rho,s,\pi^2$, by
\begin{equation}\label{cdl:eq:momentrec}
 \PP_0=1,\qquad
 \PP_{k+1}=(\rho v+sh)\PP_k+
 (v^2+\pi^2)\partial_v\PP_k+h^2\partial_h\PP_k.
\end{equation}
The first two are
\begin{align}
 \PP_1&=\rho v+sh,\label{cdl:eq:P1}\\
 \PP_2&=\rho(\rho+1)v^2+\rho\pi^2+
            2\rho s vh+s(s+1)h^2.\label{cdl:eq:P2}
\end{align}

\begin{theorem}[Complete finite moment reduction]\label{cdl:thm:moments}
Under the assumptions of Theorem~\ref{cdl:thm:addition}, every integral
obtained by inserting $\prod_jx_j^{k_j}$ in \eqref{cdl:eq:addition}
is a finite linear combination of
\begin{equation}\label{cdl:eq:momentbasis}
 \F_{R+2\ell,a}(W+d;x),\qquad \ell,d\ge0,
\end{equation}
with polynomial coefficients in $x,\pi^2,s_j,\rho_j$ and rational
factors whose denominators are products of $R+2\ell$.
The recurrence \eqref{cdl:eq:momentrec}, multiplication of the resulting
polynomials, and the reduction rules below give a terminating algorithm.
The depth shifts can then be removed by \eqref{cdl:eq:allshifts}, leaving
only finitely many spectral shifts at depth $R$.
\end{theorem}
\begin{proof}
Differentiating a bilateral transform gives
\[
 (-\partial_p)^k\{g(p)^\rho A^{-s}\}
 =g(p)^\rho A^{-s}\PP_k(v,h;\rho,s).
\]
Indeed, $-\partial_p v=v^2+\pi^2$, $-\partial_p h=h^2$, and
$-\partial_p\log(g^\rho A^{-s})=\rho v+sh$; these yield exactly
\eqref{cdl:eq:momentrec}. Weighted Schwartz bounds justify all transform
derivatives. Multiplication of the factor transforms now leaves
$g^R A^{-W}\prod_j\PP_{k_j}(v,h;\rho_j,s_j)$.

Use $v^2=g^2-\pi^2$. A monomial $v^{2q}h^k$ expands into a finite sum
of $g^{2\ell}A^{-k}$ and hence contributes
$\F_{R+2\ell,a}(W+k;x)$. For an odd power of $v$, use the contour
integration-by-parts identity
\begin{equation}\label{cdl:eq:oddreduction}
 \frac1{2\pi i}\int e^{px}v\,g(p)^D A^{-Z}\dd p
 =\frac{x}{D}\F_{D,a}(Z;x)-\frac ZD\F_{D,a}(Z+1;x).
\end{equation}
To prove it, integrate the derivative of $e^{px}g(p)^D A^{-Z}$.
The boundary terms vanish exponentially; the derivative equals
$e^{px}g^DA^{-Z}(x-Dv-Z/A)$. Since $\Rea D>0$, its denominator
$D$ is nonzero. This removes the final odd $v$ factor. The algorithm
terminates because the original moment polynomial has finite degree.
Theorem~\ref{cdl:thm:shift} gives the last reduction.
\end{proof}

At integer total $R$, equations \eqref{cdl:eq:evenbasis},
\eqref{cdl:eq:oddbasis}, and \eqref{cdl:eq:finiteLerch} reduce the result
further to finite ordinary Hurwitz--Lerch families. This conclusion does
not require any individual $\rho_j$ to be integral. Independent spectral
differentiation of the finite formula is legitimate because the orders
were never restricted to integers.

\subsection{The first moment and the complete second moment at total two}
For two factors let $R=\rho+\sigma$ and $W=s+t$. The first moment is
\begin{equation}\label{cdl:eq:firstgeneral}
 \begin{split}
 \int_\R y\F_{\rho,a}(s;y)\F_{\sigma,a}(t;x-y)\dd y
 ={}&\frac{\rho x}{R}\F_{R,a}(W;x)\\
 &+\frac{\sigma s-\rho t}{R}\F_{R,a}(W+1;x).
 \end{split}
\end{equation}
This follows from \eqref{cdl:eq:P1} and \eqref{cdl:eq:oddreduction}.
Adding the formula with the two factors exchanged gives $x$ times the
unweighted convolution, as required by $y+(x-y)=x$.

For the next formulas set $\sigma=2-\rho$, $0<\Rea\rho<2$, and
\begin{equation}\label{cdl:eq:Mdef}
 M_k(s,t)=\int_\R x^k\F_{\rho,a}(s;x)
                     \F_{\sigma,a}(t;-x)\dd x.
\end{equation}
Then
\begin{align}
 M_0(s,t)&=\E_a(W),\label{cdl:eq:M0}\\
 M_1(s,t)&=\frac{\sigma s-\rho t}{2}\E_a(W+1),\label{cdl:eq:M1}\\
 M_2(s,t)&=\frac{\pi^2\rho\sigma}{3}\E_a(W)
                  +T_{\rho,\sigma}(s,t)\E_a(W+2),\label{cdl:eq:M2}\\
 T_{\rho,\sigma}(s,t)&=
 \frac{\sigma(\sigma+1)s(s+1)+\rho(\rho+1)t(t+1)-2\rho\sigma st}{6}.
 \label{cdl:eq:T}
\end{align}
For clarity, \eqref{cdl:eq:P2} may first be written
$\rho(\rho+1)g^2-\rho^2\pi^2+2\rho s vh+s(s+1)h^2$.
Its inverse transform is
\[
 \rho(\rho+1)\F_{4,a}(W;0)-\rho^2\pi^2\E_a(W)
 -\rho s(W+1)\E_a(W+2)+s(s+1)\E_a(W+2).
\]
Substitution of \eqref{cdl:eq:C4}, followed by $\sigma=2-\rho$,
gives \eqref{cdl:eq:M2}. This derivation fixes the asymmetric order
coefficients without guessing from a symmetric special case.
At $s=t=0$ it yields the Gamma-bridge moment
\begin{equation}\label{cdl:eq:variance}
 \int_\R x^2K_\rho(x)K_{2-\rho}(-x)\dd x
 =\frac{\pi^2\rho(2-\rho)}3.
\end{equation}

\section{Continuous-split Stieltjes identities}
\label{cdl:sec:stieltjes}

Define spectral jets of the continuous-depth profile by
\begin{equation}\label{cdl:eq:Jdef}
 J_{\rho,m}^a(x)=
 \left.\partial_s^m\F_{\rho,a}(s;x)\right|_{s=0},\qquad m\ge0.
\end{equation}
These are spectral derivatives, not derivatives of the argument or the
depth. At depth one they are the order derivatives of the Lerch profiles,
and at $a=1$ of ordinary polylogarithms.

\begin{theorem}[Ordinary continuous-split Stieltjes product]
\label{cdl:thm:stieltjes}
Let $\sigma=2-\rho$, $0<\Rea\rho<2$, and $\Rea a>0$.
For $m,n\ge0$, with $M=m+n$,
\begin{equation}\label{cdl:eq:Jpair}
 \int_\R J_{\rho,m}^a(x)J_{\sigma,n}^a(-x)\dd x
 =\begin{cases}
 1,&M=0,\\
 (-1)^{M-1}M\gamma_{M-1}(a),&M\ge1.
 \end{cases}
\end{equation}
All these integrals converge absolutely. At every fixed integer total
depth $R>2$, analogous products are finite combinations of Hurwitz
spectral derivatives and Stieltjes constants obtained from
\eqref{cdl:eq:evenbasis}--\eqref{cdl:eq:oddbasis}.
\end{theorem}
\begin{proof}
Differentiate \eqref{cdl:eq:M0} independently $m$ times in $s$ and $n$
times in $t$. Theorem~\ref{cdl:thm:addition} justifies passing all
derivatives through the ordinary integral. The Laurent normalization
\eqref{cdl:eq:stieltjes} gives
\begin{equation}\label{cdl:eq:Ejets}
 \E_a^{[M]}(0)=
 \begin{cases}1,&M=0,\\(-1)^{M-1}M\gamma_{M-1}(a),&M\ge1.
 \end{cases}
\end{equation}
For larger integer total depth, differentiate the finite entire central
formulas instead. Their coefficients are finite Pochhammer polynomials;
at zero their derivatives can be written using unsigned Stirling numbers
or finite generalized harmonic sums. No numerical relation search is used.
\end{proof}

The $\rho=\sigma=1$ case is already proved in \cite{RLC}. The extension
here is that any continuous split of total depth two gives the same
Stieltjes values. These are not Hermitian inner products; no positivity
of a Stieltjes constant is implied by the displayed products.

\begin{corollary}[All first and second logarithmic Stieltjes moments]
\label{cdl:cor:Jmoments}
With the same assumptions and $M\ge1$,
\begin{equation}\label{cdl:eq:Jfirst}
 \int_\R xJ_{\rho,m}^a(x)J_{\sigma,n}^a(-x)\dd x
 =\frac{\sigma m-\rho n}{2}\E_a^{[M-1]}(1).
\end{equation}
For all $m,n\ge0$,
\begin{equation}\label{cdl:eq:Jsecond}
 \begin{split}
 &\int_\R x^2J_{\rho,m}^a(x)J_{\sigma,n}^a(-x)\dd x\\
 &\quad=\frac{\pi^2\rho\sigma}{3}\E_a^{[M]}(0)
 +\frac{\sigma(\sigma+1)m+\rho(\rho+1)n}{6}\E_a^{[M-1]}(2)\\
 &\qquad+\frac{\sigma(\sigma+1)m(m-1)+\rho(\rho+1)n(n-1)-2\rho\sigma mn}{6}
    \E_a^{[M-2]}(2).
 \end{split}
\end{equation}
Terms with negative derivative index are omitted; their polynomial
coefficients vanish identically in the relevant cases.
\end{corollary}
\begin{proof}
Apply the independent derivatives to \eqref{cdl:eq:M1} and
\eqref{cdl:eq:M2}. The coefficient in \eqref{cdl:eq:M1} has degree one;
the coefficient $T$ has degree two, so only the three displayed product
rule terms can occur. The weighted Schwartz estimates prove all required
interchanges. At $M=0$ the first moment is zero and the second moment
is \eqref{cdl:eq:variance}.
\end{proof}

A representative mixed identity, valid for every such continuous split, is
\begin{equation}\label{cdl:eq:J11}
 \begin{split}
 \int_\R x^2J_{\rho,1}^a(x)J_{\sigma,1}^a(-x)\dd x
 ={}&-\frac{2\pi^2\rho\sigma}{3}\gamma_1(a)
 +(1-\rho\sigma)\zeta(3,a)\\
 &+\left(2-\frac{2\rho\sigma}{3}\right)\zeta^{[1]}(3,a).
 \end{split}
\end{equation}
For $\rho=\sigma=1$, the $\zeta(3,a)$ term vanishes, recovering the
previous $4\zeta^{[1]}(3,a)/3-2\pi^2\gamma_1(a)/3$ formula.
The extra zeta term is necessary away from the symmetric split.

\subsection{Subtracted ordinary integrals for the spectral jets}
The spectral jets also have a real-integral construction independent of
inserting logarithms in the Fourier contour. Put
\[
 Q(u)=e^{-(a-1)u}K_\rho(x-u),\qquad q_0=K_\rho(x).
\]
For $\Rea s>-1$,
\begin{equation}\label{cdl:eq:subtracted}
 \F_{\rho,a}(s;x)=q_0+
 \frac1{\Gamma(s)}\int_0^\infty
 u^{s-1}\{Q(u)-q_0e^{-u}\}\dd u.
\end{equation}
The braces are $O(u)$ at zero and decrease exponentially at infinity.
The equality follows on $\Rea s>0$ by adding the elementary Gamma
integral and then extends to the displayed half-plane. Hence, with
\[
 I_j=\int_0^\infty\frac{Q(u)-q_0e^{-u}}u\log^j u\dd u,
\]
one obtains the ordinary convergent formulas
\begin{equation}\label{cdl:eq:Jsubtracted}
 J_{\rho,1}^a(x)=I_0,\qquad
 J_{\rho,2}^a(x)=2I_1+2\gamma I_0.
\end{equation}
More generally, if $1/\Gamma(s)=\sum_{r\ge1}c_rs^r$, then for $M\ge1$
\begin{equation}\label{cdl:eq:Jallsubtracted}
 J_{\rho,M}^a(x)=M!\sum_{j=0}^{M-1}\frac{c_{M-j}}{j!}I_j.
\end{equation}
This fixes the normalization of all zero-order jets without a finite-part
interpretation. To continue a Mellin representation farther left, one
subtracts sufficiently many Taylor terms; the entire contour formula
already supplies that continuation globally.

\section{Depth derivatives and polygamma-weighted Gamma identities}
\label{cdl:sec:depthjets}

The new depth variable is genuinely holomorphic and can now be
differentiated. For $j\ge1$ define
\begin{equation}\label{cdl:eq:Aj}
 \begin{split}
 A_j(\rho,x)={}&\mathbf1_{j=1}\log(2\pi)-\psi^{(j-1)}(\rho)\\
 &+2^{-j}\left\{
 \psi^{(j-1)}\left(\frac\rho2+\frac{ix}{2\pi}\right)+
 \psi^{(j-1)}\left(\frac\rho2-\frac{ix}{2\pi}\right)\right\}.
 \end{split}
\end{equation}
Logarithmic differentiation of \eqref{cdl:eq:K} gives
\begin{equation}\label{cdl:eq:Kdepthjets}
 \partial_\rho^qK_\rho(x)
 =K_\rho(x)\B_q(A_1(\rho,x),\ldots,A_q(\rho,x)).
\end{equation}
For real $x$ and $\Rea\rho>0$, all Gamma arguments are in the right
half-plane, so these derivatives are holomorphic and have the weighted
integrability already proved.

\begin{theorem}[Depth-jet convolution identities]\label{cdl:thm:depthjets}
For $u,v\ge0$, $R=\rho+\sigma$, and the usual positive real parts,
\begin{equation}\label{cdl:eq:depthproduct}
 \begin{split}
 &\int_\R K_\rho(y)K_\sigma(x-y)
 \B_u(A_1(\rho,y),\ldots,A_u(\rho,y))\\
 &\hspace{18mm}\cdot\B_v(A_1(\sigma,x-y),\ldots,A_v(\sigma,x-y))\dd y\\
 &\quad=K_R(x)\B_{u+v}(A_1(R,x),\ldots,A_{u+v}(R,x)).
 \end{split}
\end{equation}
For any fixed $R$ with a permitted split, every balanced derivative
satisfies, for $q\ge1$,
\begin{equation}\label{cdl:eq:balanced}
 \sum_{j=0}^q(-1)^{q-j}\binom qj
 \int_\R\partial_\rho^jK_\rho(y)
          \partial_\sigma^{q-j}K_\sigma(x-y)\dd y=0.
\end{equation}
The same balanced identity holds for spectral profiles at arbitrary
orders, and after any fixed spectral derivatives, as long as the total
spectral order is fixed during the balanced depth differentiation.
\end{theorem}
\begin{proof}
Differentiate the ordinary semigroup \eqref{cdl:eq:Ksemi} independently
in its two depths and substitute \eqref{cdl:eq:Kdepthjets}.
For \eqref{cdl:eq:balanced}, set $\sigma=R-\rho$ first and differentiate
$q$ times. The right side is then independent of $\rho$; the alternating
binomial coefficients are the product rule with
$\partial_\rho\sigma=-1$. Theorem~\ref{cdl:thm:addition} proves the
profile and spectral-jet versions with the same justification.
\end{proof}

At total depth two and $x=0$,
\[
 K_2(0)=1,\quad A_1(2,0)=\log(2\pi)-1,\quad
 A_2(2,0)=1-\frac{\pi^2}{12}.
\]
Thus the $u=1,v=0$ case can be simplified to the exact digamma average
\begin{equation}\label{cdl:eq:digammaaverage}
 \begin{split}
 &\frac12\int_\R K_\rho(x)K_{2-\rho}(-x)
 \left\{\psi\left(\frac\rho2+\frac{ix}{2\pi}\right)
       +\psi\left(\frac\rho2-\frac{ix}{2\pi}\right)\right\}\dd x\\
 &\hspace{25mm}=\psi(\rho)-1.
 \end{split}
\end{equation}
At total depth one, the useful anchors are
$K_1(0)=1/2$, $A_1(1,0)=\log(\pi/2)$, and
$A_2(1,0)=\pi^2/12$.

An unbalanced depth derivative at total depth two generally involves
$\partial_R\F_{R,a}(s;x)|_{R=2}$, not a derivative of $\E_a(s)$ with
respect to a nonexistent depth parameter. The distinction matters:
Section~\ref{cdl:sec:obstruction} shows why logarithmic depth multipliers
cannot simply be declared part of the old finite meromorphic basis.

\section{Shift derivatives and all anchored primitive orders}
\label{cdl:sec:primitives}

\begin{theorem}[Shift calculus and removable primitive resonances]
\label{cdl:thm:primitives}
For every integer $k\ge0$,
\begin{equation}\label{cdl:eq:aderivative}
 \partial_a^k\F_{\rho,a}(s;x)=(-1)^k(s)_k\F_{\rho,a}(s+k;x).
\end{equation}
Fix a base point $b$ in the right half-plane. For $r\ge1$, let
$\mathcal J_b^r$ be the $r$-fold ordinary primitive in $a$, with all
first $r$ initial data zero at $a=b$. Then
\begin{equation}\label{cdl:eq:aprimitive}
 \begin{split}
 \mathcal J_b^r\F_{\rho,a}(s;x)
 =\frac{(-1)^r}{(s-r)_r}\biggl[&\F_{\rho,a}(s-r;x)\\
 &-\sum_{k=0}^{r-1}\frac{(-1)^k(s-r)_k(a-b)^k}{k!}
             \F_{\rho,b}(s-r+k;x)\biggr].
 \end{split}
\end{equation}
All apparent singularities at $s=1,\ldots,r$ are removable. The formula
is entire in $s$ after those values are filled in.
\end{theorem}
\begin{proof}
Differentiating $(p+a-1)^{-s}$ under the exponentially convergent
contour proves \eqref{cdl:eq:aderivative}. Away from the displayed
resonances, an $r$th antiderivative is
$(-1)^r\F_{\rho,a}(s-r;x)/(s-r)_r$.
Subtract its Taylor polynomial of degree $r-1$ at $b$ using
\eqref{cdl:eq:aderivative}; this is exactly \eqref{cdl:eq:aprimitive}.
It has the required derivative and the required initial conditions.

The actual anchored primitive also has the finite-path integral
\[
 \frac1{(r-1)!}\int_b^a(a-u)^{r-1}\F_{\rho,u}(s;x)\dd u,
\]
where the segment stays in the right half-plane. It is entire in $s$,
so equality away from a finite set proves removability there. Equivalently,
at $s\in\{1,\ldots,r\}$ the negative-order expression
\eqref{cdl:eq:negative} is a polynomial in $a$ of degree at most $r-1$;
its Taylor remainder in the numerator vanishes. The finite-path formula
is the normalization, not an unspecified additive constant.
\end{proof}

For $r=1$ the formula reads
\begin{equation}\label{cdl:eq:firstaprimitive}
 \int_b^a\F_{\rho,u}(s;x)\dd u
 =\frac{\F_{\rho,b}(s-1;x)-\F_{\rho,a}(s-1;x)}{s-1}.
\end{equation}
At $s=1$ take the spectral derivative of the numerator; assigning zero
to each of its two equal zeroth-order terms before dividing is invalid.

There is also an all-order primitive in the argument after the explicit
covariant weight is inserted. Let
$G_s(x)=e^{(a-1)x}\F_{\rho,a}(s;x)$. The ladder gives
$G_{s+1}'=G_s$. Therefore
\begin{equation}\label{cdl:eq:xprimitive}
 \mathcal J_{x_0}^r G_s(x)=G_{s+r}(x)
 -\sum_{j=0}^{r-1}\frac{(x-x_0)^j}{j!}G_{s+r-j}(x_0).
\end{equation}
For $a=1$ this is an ordinary unweighted argument primitive of
$\F_{\rho,1}$, including the standard polylogarithm ladder at depth one.
For other $a$, the weight is part of the statement and must not be omitted.

\subsection{Gamma primitives of the Stieltjes product family}
At total depth two, Hurwitz differentiation gives
\begin{equation}\label{cdl:eq:Eprimitive}
 \int_b^a\E_u(s)\dd u=\zeta(s,b)-\zeta(s,a).
\end{equation}
Both sides are entire in $s$; the two residues at $s=1$ cancel.
At $s=0$ this is $a-b$. Differentiating at $s=0$, and using
$\zeta^{[1]}(0,a)=\log\Gamma(a)-\tfrac12\log(2\pi)$, gives
\begin{equation}\label{cdl:eq:gammaprimitive}
 \int_b^a\gamma_0(u)\dd u=\log\Gamma(b)-\log\Gamma(a).
\end{equation}
The logarithm of Gamma is the analytic branch on the right half-plane.
Combining with Theorem~\ref{cdl:thm:stieltjes}, for $m+n=1$,
\begin{equation}\label{cdl:eq:tripleprimitive}
 \int_b^a\int_\R J_{\rho,m}^u(x)J_{2-\rho,n}^u(-x)\dd x\dd u
 =\log\Gamma(b)-\log\Gamma(a).
\end{equation}
Uniform weighted bounds on the compact segment justify the interchange.
For arbitrary $M=m+n\ge1$, the same procedure gives
\begin{equation}\label{cdl:eq:allprimitivejets}
 \int_b^a\int_\R J_{\rho,m}^u(x)J_{2-\rho,n}^u(-x)\dd x\dd u
 =\zeta^{[M]}(0,b)-\zeta^{[M]}(0,a).
\end{equation}

For completeness, every derivative in the Stieltjes parameter has a
finite spectral-zeta expression. If $k,M\ge1$, then
\begin{equation}\label{cdl:eq:stieltjesaderivatives}
 \begin{split}
 (-1)^{M-1}M\,\partial_a^k\gamma_{M-1}(a)
 =(-1)^k\sum_{j=1}^{\min(M,k+1)}
 \frac{M!}{(M-j)!}\stir{k+1}{j}\zeta^{[M-j]}(k+1,a).
 \end{split}
\end{equation}
Indeed, $\partial_a^k\E_a(s)=(-1)^k(s)_{k+1}\zeta(s+k+1,a)$.
Expand $(s)_{k+1}=\sum_j\stir{k+1}{j}s^j$ and differentiate at zero.
This is a classical Hurwitz-calculus consequence included to make the
derivative--primitive closure explicit; it is not advertised as a new
formula for the Stieltjes constants themselves.

\section{A precise obstruction to finite integer-depth reduction}
\label{cdl:sec:obstruction}

It is useful to state the obstruction in a form that can actually be
proved. A statement that ``no finite special-function expression exists''
without specifying the allowed class would be unjustified.

\begin{definition}[Local integer-depth transform class]
\label{cdl:def:class}
Fix a punctured neighborhood of $p=1$. Let $\mathscr A$ be the class
obtained by finite sums, products, and $p$-derivatives from multipliers
\[
 g(p)^r(p+a_j-1)^{-u_j},\qquad r\in\N,\quad \Rea a_j>0,\quad u_j\in\C,
\]
allowing multiplication by single-valued meromorphic coefficients near
$p=1$ and any fixed spectral derivatives in the $u_j$. Every logarithm
of $p+a_j-1$ is chosen on the full sufficiently small disk, where its
argument is nonzero. Thus every member of $\mathscr A$ is single-valued
meromorphic near $p=1$.
\end{definition}

This class contains the transform operations used for finite integer-depth
spectral shifts, polynomial argument moments, and their order derivatives.
It is intentionally not the class of every possible special function.

\begin{theorem}[Monodromy obstruction]\label{cdl:thm:obstruction}
Let $\Rea a>0$ and $s\in\C$. If $\Rea\rho>0$ and $\rho\notin\N$, then
\begin{equation}\label{cdl:eq:obstructed}
 g(p)^\rho(p+a-1)^{-s}\notin\mathscr A.
\end{equation}
At an integer $r\ge1$, every positive depth derivative
\begin{equation}\label{cdl:eq:depthobstruction}
 g(p)^r(p+a-1)^{-s}(\log g(p))^q,\qquad q\ge1,
\end{equation}
also lies outside this local meromorphic class.
\end{theorem}
\begin{proof}
Near $p=1$, the function $g$ has a simple pole and a nonvanishing
holomorphic leading factor. Continuing its logarithm once around a
positively oriented small loop changes $\log g$ by $-2\pi i$.
The factor $(p+a-1)^{-s}$ is holomorphic and nonzero on that disk,
because $a\ne0$. Hence the multiplier in \eqref{cdl:eq:obstructed}
acquires the monodromy factor $e^{-2\pi i\rho}$, which is nontrivial
precisely when $\rho$ is not an integer. This contradicts the
single-valuedness of every member of $\mathscr A$.

For \eqref{cdl:eq:depthobstruction}, the nonzero meromorphic factor
$g^r(p+a-1)^{-s}$ is unchanged, but $(\log g)^q$ changes to
$(\log g-2\pi i)^q$. Their difference is a nonzero polynomial in
$\log g$ of degree $q-1$, so it cannot vanish identically.
An equality on the original transform strip would force equality of
these analytic continuations by uniqueness; the local contradiction
therefore excludes such an equality there as well.
\end{proof}

This proves why the finite recurrence reduces integer depths to the two
elementary seeds, while a noninteger residue class of depths retains
its own Gamma kernel. The finite shifts
$\rho,\rho+2,\ldots,\rho+2n$ still reduce to one continuous-depth seed
by Theorem~\ref{cdl:thm:shift}. There is no contradiction between that
finite reduction and the obstruction to changing a noninteger seed
into the old integer-depth class.

The theorem is about identities of transform families. It proves neither
arithmetic independence of isolated central values nor the impossibility
of representations in larger classes with nontrivial monodromy. In
particular it does not obstruct the exact odd-order values of
Section~\ref{cdl:sec:reflection}, or a finite closure after several
noninteger depths sum to an integer.

\section{Source audit, verification, and integration status}
\label{cdl:sec:audit}

\subsection{Precisely what was inspected}
Current repository inspection was pinned to \pin. It covered the
canonical manuscript README, the full chapter
\src{chapters/05-logistic-resolvent-identities.tex}, directory/tree
metadata for the manuscript and incoming area, and the relevant visible
intake rules in \src{docs/incoming/README.md}. The substantial prior
mathematical reading was of the accessible Library TEX copy of
\emph{Reciprocal Lerch Convolutions}: its transform and entire-order
proofs, integer-depth central-factorial formulas, Stieltjes and harmonic
coefficient sections, and the noninteger-depth research question.
That older report cites its own source pin
\src{0c740647c48baec7f0b4192723e977d69917dd52}; the two pins must not
be conflated. Introductory snippets of the later commensurate-Lerch
report were sufficient to avoid claiming its scale/unequal-center agenda
as the present target, but its full proofs were not audited here.

The binary incoming ZIP packages were listed, not read in full. The
Library source is not asserted to be byte-identical to any currently
listed archive. This is a targeted research continuation, not an audit
of every chapter, every incoming report, or every existing conjecture.

\subsection{Corrections and safeguards}
\textbf{No incorrect source formula was established in the inspected
proof sections.} Accordingly this package does not label an unverified
suspicion as an erratum. It does identify several necessary safeguards
for extending those results:

\begin{enumerate}[label=(\alph*)]
\item A squared-modulus Gamma formula is not holomorphic in complex depth.
Use the product \eqref{cdl:eq:K}, which agrees with the squared modulus
only for real positive depth.
\item Negative spectral orders require \eqref{cdl:eq:negative}, not
substitution into a divergent positive-order Mellin integral.
The value $\F_{2,a}(0;0)=1$ is the simplest normalization test.
\item Pole--zero combinations in \eqref{cdl:eq:reflection}, the finite
Hurwitz basis, and the anchored primitive formula must be combined
before a resonant parameter is substituted.
\item Integer-depth finite formulas cannot be differentiated in their
integer label. Depth differentiation is justified here only after
constructing the holomorphic Gamma family. The resulting logarithmic
multipliers generally leave the old meromorphic class.
\item ``Finite closure'' must name its function class. Theorem
\ref{cdl:thm:obstruction} supplies such a class and explicitly excludes
any inference about independence of isolated constants.
\item The symmetric split $\rho=\sigma=1$ hides the additional
$(1-\rho\sigma)\zeta(3,a)$ term in \eqref{cdl:eq:J11}. Keeping it is
necessary when extending the earlier moment identity to unequal depths.
\end{enumerate}

The reference chapter's beta-density generating parameter and the present
reciprocal-convolution depth are different constructions. Their common
use of beta integrals is not evidence that a conjecture in one geometry
has been proved in the other. In particular the Gaussian $S_6$ and $S_8$
status is unchanged by this work.

\subsection{Executed checks and their evidentiary meaning}
The package includes executable exact and numerical verification, with
machine-readable results. The exact program performs 241 finite
assertions and three negative controls. The principal families are
summarized below.
\begin{center}
\begin{tabular}{@{}p{0.79\linewidth}r@{}}
\toprule
Exact algebra family & Checks\\\midrule
Two independent Gamma-coefficient recurrences and homogeneous weights &56\\
Finite products versus Bell--harmonic derivatives &108\\
Depth-shift product normalizations and recurrences &13\\
Transform moment recurrence, Bell comparison, and leading coefficients &27\\
First/second moment coefficient identities and explicit harmonic examples &5\\
Shifted generalized-harmonic sums &32\\\midrule
Total exact assertions &241\\\bottomrule
\end{tabular}
\end{center}
The coefficient table covers $1\le m\le4$ and $0\le q\le6$.
Formal zeta symbols are used in a polynomial ring for exact algebra;
this does not assume their actual arithmetic independence.

The numerical suite executed 88 diagnostics at 50 working decimal digits:
38 kernel, transform, and moment tests; 28 centered, reflection, harmonic,
and polylogarithm tests; and 22 Stieltjes, depth-jet, and primitive tests. Every test used the
scaled residual $|L-R|/(1+|R|)$ with acceptance threshold $10^{-30}$;
the largest observed scaled residual was approximately $3.54\times10^{-47}$.
The tested cases include complex noninteger depths, arbitrary complex
orders, negative spectral orders, shifts below $1/2$ requiring a
nonsymmetric contour, direct Gamma convolutions, the full harmonic
parameter identity, and a genuine Gamma-weighted trilogarithm integral.
The subtracted Mellin jet calculation is independent of contour logarithm
insertion. Alternating sums use mpmath's convergence acceleration;
they are not finite tail enclosures.

\textbf{Working precision and small residuals are not rigorous error
bounds.} These diagnostics can expose mistakes but do not prove the
infinite-parameter identities. The analytic proofs establish those
identities. No interval certification, proof-assistant kernel check,
independent refereeing, or exhaustive literature-priority verification
is claimed. The receipts record the actual package versions and each
individual comparison.

\subsection{Additive integration recommendation}
This is an addition answering question 10 of the reciprocal-Lerch report,
not a replacement for the canonical manuscript. A natural relative
placement beneath that continuation is
\src{continuations/continuous-depth-lerch/}; the maintainer should select
the actual destination on the Mellin--Stieltjes spine in accordance with
the intake README. This proposed path is not asserted to exist already.
All theorem and equation labels in this article use the prefix
\src{cdl:}. The accompanying integration note maps the proved statements,
classical recoveries, and remaining questions. No repository or persistent
Library mutation was performed.

\section{Further research questions and concrete next targets}
\label{cdl:sec:research}

The following questions extend the exact-identity programme rather than
asking for better inequalities. The distinction between a proved formula
and an open proposed extension is maintained in every item.

\begin{enumerate}[label=Q\arabic*.]
\item \textbf{Depth zero and negative depth.}
Continue the normalized semigroup through $\rho=0$ as distributions,
identify the delta contribution, and derive the explicit residue/contact
terms needed for negative real parts of $\rho$. The positive-half-plane
kernel alone does not justify dropping Gamma poles that cross an
argument contour. The left-half-plane Gamma-transform work in
\cite{CS} supplies a classical comparison, not an automatic answer for
the coupled resolvent.

\item \textbf{Rational base-depth harmonic masters.}
Replace $2+\epsilon$ in Theorem~\ref{cdl:thm:harmonicmaster} by
$2a+\epsilon$ with rational $a>0$. Derive the corresponding shifted
harmonic Bell combinations and classify the residue-class combinations
whose Gamma coefficients reduce to ordinary zeta products, rather than
general cyclotomic constants. Every proposed reduction should be an
identity of the full generating functions.

\item \textbf{Even polylogarithm orders under Gamma weights.}
The odd family \eqref{cdl:eq:polylogGamma} is explicit because of the
polynomials $U_m$. Determine an equally useful exact description for
even orders. The reflected values $Z_\rho^{[1]}(1-2m)$ in
\eqref{cdl:eq:evenreflection} identify the missing spectral data; they
must not be confused with the known trivial zeros themselves.

\item \textbf{Independent resolvent shifts in the centered reflection.}
For a product $g(p)^R\prod_j(p+a_j-1)^{-s_j}$, construct a centered
reflection identity that retains all independent spectral jets and
states the branch geometry at shift collisions. Existing unequal-center
expansions should be used as input rather than re-labeled as new;
the specific target is an exact reflected function and its collision law.

\item \textbf{Unbalanced depth jets at integer total depth.}
Find a canonical enlarged basis containing
$g(p)^R(\log g(p))^q(p+a-1)^{-s}$ and derive its special-value identities.
Theorem~\ref{cdl:thm:obstruction} rules out only the old local meromorphic
class. A logarithmically extended Lerch family may have useful finite
recurrences even though it is outside that class.

\item \textbf{A multivariate tilt generating function.}
The finite moment recurrence determines every Taylor coefficient of
convolutions with factors $e^{\lambda_jx_j}$. Obtain a compact analytic
function of all $\lambda_j$, with explicit confluence and branch rules.
It must reproduce \eqref{cdl:eq:firstgeneral} and \eqref{cdl:eq:M2}
without losing the asymmetric order terms.

\item \textbf{Products of polylogarithms under one Gamma weight.}
Evaluate integrals containing several independent polylogarithms in the
same variable against $G_\rho(x)$. Their transform geometry is not the
simple reciprocal convolution treated here. Determine which balanced
odd-part combinations admit Gamma-quotient generating functions, and
which require multiple Hurwitz--Lerch functions.

\item \textbf{Spectral derivatives on the odd evaluation lattice.}
Formula \eqref{cdl:eq:oddF} evaluates the entire function at odd integers
but does not determine its order derivatives there. Use the reflection
formula and ordinary contour derivatives to derive those derivatives,
including their Gamma, Stieltjes, and $Z_\rho$-jet terms. Differentiating
an identity known only for discrete $m$ is not a valid shortcut.

\item \textbf{Canonical reduction of continuous moment polynomials.}
The moment algorithm terminates, but its output is not a unique normal
form. Determine the exact module of relations among shifted
$\F_{R,a}(W+d;x)$ induced by the depth recurrence and covariant ladder.
A canonical symbolic reduction should be proved without assuming
independence of special values.

\item \textbf{A genuine connection to the Gaussian harmonic candidates.}
Investigate whether the Gamma-weighted odd-polylogarithm family can
produce the missing mixed-color functional relations relevant to the
canonical $S_6$ and $S_8$ problems. Any claimed advance must exhibit the
new convergent relation or certificate explicitly. The continuous-depth
semigroup by itself does not supply such a relation.

\item \textbf{Formalization and literature reconciliation.}
Formalize the polynomial recurrences, Bell coefficients, moment
reductions, and primitive Taylor remainders as a finite algebra layer.
Isolate weighted Schwartz holomorphy, Fourier inversion, and Mellin
subtraction as separate analytic obligations. Independently compare the
precise harmonic master and its coefficient families with the classical
Euler-sum and generalized-Lerch literature before assigning priority
labels to individual identities.
\end{enumerate}

\section*{Conclusion}
\addcontentsline{toc}{section}{Conclusion}
The specified noninteger-depth question has an explicit positive-half-plane
answer: a Gamma kernel with a holomorphic entire-order resolvent extension.
Depth and spectral order add independently. Integer \emph{total} depths
retain finite zeta closure, all polynomial moments admit a terminating
calculus, and continuous splits give ordinary Stieltjes products.
The centered reflection exposes two exact odd-order families: one with
Gamma--Hurwitz coefficients and one polynomial in depth. Their combination
produces an all-order harmonic generating identity and a continuous
Gamma-weighted polylogarithm evaluation family. The finite-closure
obstruction is local and precisely stated, while the unresolved boundary,
multishift, and spectral-jet extensions remain explicit research targets.

\appendix
\section{Selected exact harmonic coefficient values}
\label{cdl:app:coeffs}
Let $S_{m,q}$ denote the left side of \eqref{cdl:eq:mastercoeff}. The
following entries are generated by the exact program. Powers of $\pi$
here are simply the conventional simplifications of even zeta values;
the complete formal-zeta table is retained in the JSON artifact.
\begin{center}
\renewcommand{\arraystretch}{1.35}
\begin{tabular}{@{}rrl@{}}
\toprule
$m$ & $q$ & $S_{m,q}$\\\midrule
% BEGIN GENERATED HARMONIC TABLE
1 & 0 & $\frac{\pi^{2}}{12}$ \\
1 & 1 & $- \frac{\zeta\left(3\right)}{2}$ \\
1 & 2 & $\frac{\pi^{4}}{720}$ \\
1 & 3 & $- \frac{3 \zeta\left(5\right)}{2} + \frac{\pi^{2} \zeta\left(3\right)}{4}$ \\
1 & 4 & $- 3 \zeta^{2}\left(3\right) - \frac{67 \pi^{6}}{20160}$ \\
2 & 0 & $\frac{7 \pi^{4}}{720}$ \\
2 & 1 & $- \frac{\pi^{2} \zeta\left(3\right)}{12} - \frac{\zeta\left(5\right)}{2}$ \\
2 & 2 & $\frac{\zeta^{2}\left(3\right)}{2} + \frac{23 \pi^{6}}{12096}$ \\
2 & 3 & $- \frac{15 \zeta\left(7\right)}{4} - \frac{\pi^{2} \zeta\left(5\right)}{8} + \frac{\pi^{4} \zeta\left(3\right)}{96}$ \\
2 & 4 & $- \frac{3 \pi^{2} \zeta^{2}\left(3\right)}{4} + 3 \zeta\left(3\right) \zeta\left(5\right) + \frac{43 \pi^{8}}{57600}$ \\
% END GENERATED HARMONIC TABLE
\bottomrule
\end{tabular}
\end{center}
The table is not the proof of the infinite family. Theorem
\ref{cdl:thm:harmonicmaster} supplies that proof for arbitrary $m,q$.

\begin{thebibliography}{99}
\addcontentsline{toc}{section}{References}
\bibitem{Repo}
V. Reshetnikov, \emph{ProveIt: canonical polylogarithm manuscript},
repository snapshot \pin. Inspected source paths include
\src{Analysis/Polylogarithms/docs/manuscript/README.md} and
\src{Analysis/Polylogarithms/docs/manuscript/chapters/05-logistic-resolvent-identities.tex}.
\url{https://github.com/VladimirReshetnikov/ProveIt}.

\bibitem{RLC}
\emph{Reciprocal Lerch Convolutions}, research continuation prepared for the
ProveIt programme, 10 October 2026. Accessible Library TEX source
\src{ProveIt_Reciprocal_Lerch_Convolutions_2026-10-10.tex}, especially
Sections 2--4, 6, and Section 12, question 10. The report's own repository
source pin is \src{0c740647c48baec7f0b4192723e977d69917dd52}.

\bibitem{CS}
D. Chakrabarti and G. K. Srinivasan, \emph{On a remarkable formula of
Ramanujan}, 2012, arXiv:1206.5122. Theorem 1.1 gives the Gamma-product
Fourier transform for complex parameters in the right half-plane.
\url{https://arxiv.org/abs/1206.5122}.

\bibitem{Schoutens}
W. Schoutens, \emph{The Meixner Process: Theory and Applications in
Finance}, EURANDOM Report 2002-004, 12 February 2002, Section 2,
printed pp. 5--6, for the Gamma density and its convolution parameter.
\url{https://www.eurandom.tue.nl/reports/2002/004-report.pdf}.

\bibitem{DLMFBeta}
NIST Digital Library of Mathematical Functions,
Section 5.12, \emph{Beta Function}, Euler beta integrals.
\url{https://dlmf.nist.gov/5.12}.

\bibitem{DLMFGamma}
NIST Digital Library of Mathematical Functions,
Sections 5.7 and 5.15, \emph{Series Expansions} and
\emph{Polygamma Functions}.
\url{https://dlmf.nist.gov/5.7}; \url{https://dlmf.nist.gov/5.15}.

\bibitem{DLMFLerch}
NIST Digital Library of Mathematical Functions,
Section 25.14, \emph{Lerch's Transcendent}.
\url{https://dlmf.nist.gov/25.14}.

\bibitem{DLMFZeta}
NIST Digital Library of Mathematical Functions,
Section 25.11, \emph{Hurwitz Zeta Function}, including parameter
recurrence, derivatives, and the log-Gamma special value.
\url{https://dlmf.nist.gov/25.11}.
\end{thebibliography}
\end{document}
